\documentclass[letterpaper,10pt]{article}

\usepackage{latexsym}
\usepackage[empty]{fullpage}
\usepackage{titlesec}
\usepackage{marvosym}
\usepackage[usenames,dvipsnames]{color}
\usepackage{verbatim}
\usepackage{enumitem}
%\usepackage{xcolor}
\usepackage{color}% http://ctan.org/pkg/color

\usepackage[pdftex,colorlinks=false,pdfborderstyle={/S/U/W 1},urlbordercolor={0 0 0}]{hyperref}
\usepackage{fancyhdr}
\usepackage[italian,ngerman,english]{babel}
\usepackage{tabularx}
% only for pdflatex
% \input{glyphtounicode}

% fontawesome
\usepackage{fontawesome5}

% fixed width
\usepackage[scale=0.90,lf]{FiraMono}

% light-grey
\definecolor{light-grey}{gray}{0.83}
\definecolor{dark-grey}{gray}{0.3}
\definecolor{text-grey}{gray}{.08}

\DeclareRobustCommand{\ebseries}{\fontseries{eb}\selectfont}
\DeclareTextFontCommand{\texteb}{\ebseries}

% custom underilne
\usepackage{contour}
\usepackage[normalem]{ulem}
\renewcommand{\ULdepth}{1.8pt}
\contourlength{0.8pt}
\newcommand{\myuline}[1]{%
  \uline{\phantom{#1}}%
  \llap{\contour{white}{#1}}%
}


% font: arial
\usepackage{helvet}
\renewcommand{\familydefault}{\sfdefault}
% font: helvetica-style
%\usepackage{tgheros}
%\renewcommand*\familydefault{\sfdefault} 
%% Only if the base font of the document is to be sans serif
\usepackage[T1]{fontenc}


\pagestyle{fancy}
\fancyhf{} % clear all header and footer fields
\fancyfoot{}
\renewcommand{\headrulewidth}{0pt}
\renewcommand{\footrulewidth}{0pt}

% Adjust margins
\addtolength{\oddsidemargin}{-0.5in}
\addtolength{\evensidemargin}{0in}
\addtolength{\textwidth}{1in}
\addtolength{\topmargin}{-.8in}
\addtolength{\textheight}{1.5in}

\urlstyle{same}

\raggedbottom
\raggedright
\setlength{\tabcolsep}{0in}

% Sections formatting - serif
% \titleformat{\section}{
%   \vspace{2pt} \scshape \raggedright\large % header section
% }{}{0em}{}[\color{black} \titlerule \vspace{-5pt}]

% TODO EBSERIES
% sans serif sections
\titleformat {\section}{
    \bfseries \vspace{2pt} \raggedright \large % header section
}{}{0em}{}[\color{light-grey} {\titlerule[2pt]} \vspace{-4pt}]

% only for pdflatex
% Ensure that generate pdf is machine readable/ATS parsable
% \pdfgentounicode=1

%-------------------------
% Custom commands
\newcommand{\resumeItem}[1]{
  \item\small{
    {#1 \vspace{-1pt}}
  }
}

\newcommand{\resumeSubheading}[4]{
  \vspace{-1pt}\item
    \begin{tabular*}{\textwidth}[t]{l@{\extracolsep{\fill}}r}
      \textbf{#1} & {\color{dark-grey}\small #2}\vspace{1pt}\\ % top row of resume entry
      \textit{#3} & {\color{dark-grey} \small #4}\\ % second row of resume entry
    \end{tabular*}\vspace{-4pt}
}

\newcommand{\resumeSubSubheading}[2]{
    \item
    \begin{tabular*}{\textwidth}{l@{\extracolsep{\fill}}r}
      \textit{\small#1} & \textit{\small #2} \\
    \end{tabular*}\vspace{-7pt}
}

\newcommand{\resumeProjectHeading}[2]{
    \item
    \begin{tabular*}{\textwidth}{l@{\extracolsep{\fill}}r}
      #1 & {\color{dark-grey}} \\
    \end{tabular*}\vspace{-4pt}
}

\newcommand{\resumeSubItem}[1]{\resumeItem{#1}\vspace{-4pt}}

\renewcommand\labelitemii{$\vcenter{\hbox{\tiny$\bullet$}}$}

% CHANGED default leftmargin  0.15 in
\newcommand{\resumeSubHeadingListStart}{\begin{itemize}[leftmargin=0in, label={}]}
\newcommand{\resumeSubHeadingListEnd}{\end{itemize}}
\newcommand{\resumeItemListStart}{\begin{itemize}}
\newcommand{\resumeItemListEnd}{\end{itemize}}

\color{text-grey}

%-------------------------------------------
%%%%%%  RESUME STARTS HERE  %%%%%%%%%%%%%%%%%%%%%%%%%%%%


\begin{document}
\selectlanguage{italian}
%----------HEADING----------
% See icons: http://mirrors.ibiblio.org/CTAN/fonts/fontawesome/doc/fontawesome.pdf
\begin{center}
    \textbf{\Huge Niccol\`o Zuppichini} \\ \vspace{5pt}
    \small \faPhone* \texttt{+41 76 275 07 06} \hspace{1pt} $|$
    \hspace{1pt} \faEnvelope \hspace{2pt} \texttt{niccolo.zuppichini@gmail.com} \hspace{1pt}  \\
    \hspace{1pt} \faLinkedin \hspace{2pt} \texttt{\href{https://www.linkedin.com/in/niccolo-zuppichini}{www.linkedin.com/in/niccolo-zuppichini}} \hspace{1pt} \\
%    \hspace{1pt} \faGithub \hspace{2pt} \texttt{rage997} \hspace{1pt} $|$ \\
    \hspace{1pt} \faHome \hspace{2pt}\texttt{\href{niccolozuppichini.com}{niccolozuppichini.com}}   \hspace{1pt} $|$
    \hspace{1pt} \faMapMarker* \hspace{2pt}\texttt{Bellinzona, Svizzera}
    \\ \vspace{-3pt}
\end{center}

%-----------EXPERIENCE-----------
\section{ESPERIENZA}
  \resumeSubHeadingListStart
  
%-----------Lemonate-----------
\resumeSubheading
    {\href{https://lemonate.io/}{Lemonate}}{aprile 2025 -- oggi}
    {Senior Software Engineer}{Schaffhausen, CH (Remoto)}
    \resumeItemListStart
      \resumeItem{Contribuisco allo sviluppo del core della piattaforma 3D interattiva di Lemonate, abilitando rendering in tempo reale e deployment di applicazioni browser-based complesse.}
      \resumeItem{Sviluppo funzionalit\`a del motore di rendering, inclusi sistemi di particelle in tempo reale, pipeline di post-processing e architettura grafica a basso livello.}
      \resumeItem{Implemento tecniche avanzate di rendering e ottimizzazioni delle prestazioni in un motore basato su WebGL per visualizzazioni 3D multipiattaforma.}
      \resumeItem{Guido le decisioni architetturali del motore e sviluppo tooling per snellire i flussi di lavoro 3D negli ambienti web.}
    \resumeItemListEnd

%-----------Mindcraft AG-----------
\resumeSubheading
    {\href{https://mndcr.com/}{Mindcraft AG}}{ottobre 2022 -- aprile 2025}
    {Senior Software Engineer}{Berna, CH (Remoto)}
    \resumeItemListStart
      \resumeItem{Progetto e implemento un motore di rendering web ad alte prestazioni per la piattaforma Meshvalley.}
      \resumeItem{Creo pipeline di scansione 3D end-to-end, inclusa l'elaborazione degli asset, l'ottimizzazione e la preparazione per il rendering in tempo reale.}
      \resumeItem{Architetto sistemi per l'hosting, lo streaming e la visualizzazione interattiva di dataset 3D di grandi dimensioni nel browser.}
      \resumeItem{Sviluppo algoritmi geometrici e flussi di lavoro di elaborazione dati per gestire in modo efficiente ambienti scansionati complessi.}
    \resumeItemListEnd

%-----------Panua-----------
    \resumeSubheading
      {\href{https://panua.ch/}{Panua}}{luglio 2022 -- novembre 2023}
      {Web Developer / DevOps Engineer (Freelance)}{Lugano, CH (Remoto)}
      \resumeItemListStart
      \resumeItem{Collaboro con il team di Panua su base oraria, fornendo competenze in sviluppo web, DevOps e ingegneria del software.}
      \resumeItem{Implemento un software per generare, validare e autenticare le impronte digitali delle macchine degli utenti, prevenendo la pirateria dei prodotti software di Panua Technologies e rafforzando la sicurezza dei prodotti aziendali.}
      \resumeItem{Garantisco alta disponibilit\`a e prestazioni delle piattaforme software di Panua.}
      \resumeItem{Sviluppo pipeline automatizzate di deployment e testing con Docker, migliorando l'efficienza e l'affidabilit\`a dei rilasci software.}
    \resumeItemListEnd

%-----------HEGIAS-----------
    \resumeSubheading
      {\href{https://hegias.com/en/}{Hegias}}{settembre 2018 -- settembre 2020}
      {Software Engineer | DevOps Manager}{Lugano \& Zurigo, CH (Ibrido)}
      \resumeItemListStart
       \resumeItem{Guido un team Agile di 5 sviluppatori verso il deployment continuo, aumentando la produttivit\`a del 20\% grazie al pair programming.}
        \resumeItem{Responsabile dell'infrastruttura di server Unix in grado di gestire le richieste di migliaia di clienti.}
        \resumeItem{Risolvo in autonomia i problemi di prestazioni del motore di rendering implementando tecniche avanzate di rendering, compressione di oggetti 3D e classificazione di oggetti 3D tramite machine learning.}
        \resumeItem{Riduco i tempi di esecuzione della pipeline di elaborazione degli oggetti 3D di oltre il 250\% distribuendo l'infrastruttura server su nodi diversi.}
        \resumeItem{L'azienda ha raccolto oltre 2,1 milioni di franchi svizzeri di finanziamenti ed \`e stata selezionata tra le \href{https://www.eu-startups.com/2019/02/10-swiss-startups-to-watch-in-2019/}{10 startup da tenere d'occhio}.}
%         \resumeItem{Tecnologie utilizzate: Unity, Numpy, PyMesh, Node.js, MongoDB, three.js, WebGL}
      \resumeItemListEnd
      
%-----------CS-----------
    \resumeSubheading
      {\href{https://www.ics.usi.ch/}{Istituto di Scienza Computazionale}}{giugno 2018 -- settembre 2018}
      {Ricercatore Junior}{Lugano, CH}
      \resumeItemListStart
        \resumeItem{Contribuisco allo sviluppo di \href{https://bitbucket.org/zulianp/utopia/src/master/}{Utopia}, una libreria di calcolo parallelo C++ all'avanguardia per la risoluzione di problemi non lineari ad alta dimensionalit\`a.}
        \resumeItem{Il codice sorgente completo \`e disponibile sul \href{https://bitbucket.org/zulianp/utopia/src/master/}{repository del maintainer del progetto}}
    \resumeItemListEnd

  \resumeSubHeadingListEnd


% --------- Education -----------
\section {ISTRUZIONE}
\resumeSubHeadingListStart
   \resumeSubheading
     {Universit\`a della Svizzera Italiana}{settembre 2019 -- giugno 2022}
     {Laurea Magistrale in Scienze Computazionali}{Lugano, CH}
       \resumeItemListStart
       \resumeItem 
           {\textbf{Tesi}: \href{https://github.com/Rage997/MasterThesis}{Modellazione delle specie invasive tramite modelli di eventi relazionali.}}
     \resumeItem {\textbf{Materie}: Ottimizzazione Stocastica, Analisi Numerica, EDP/ODE, Deep Learning.}
     \resumeItem{\textbf{Media dei voti}: 7.66}
       \resumeItemListEnd
 \resumeSubHeadingListEnd
\resumeSubHeadingListStart
   \resumeSubheading
     {Universit\`a della Svizzera Italiana}{settembre 2016 -- giugno 2019}
     {Laurea Triennale in Informatica}{Lugano, CH}
       \resumeItemListStart
       \resumeItem 
           {\textbf{Progetto}: \href{https://github.com/Rage997/BachelorProject}{Una pipeline robusta per l'elaborazione di oggetti 3D.}}
     \resumeItem {\textbf{Materie}: Strutture Dati, Algoritmi, Database, Sistemi di Calcolo, Machine Learning.}
      \resumeItem{\textbf{Media dei voti}: 7.45}
       \resumeItemListEnd
 \resumeSubHeadingListEnd


%-----------PROJECTS-----------

\section{PROGETTI}
    \resumeSubHeadingListStart
     \resumeProjectHeading
         {\textbf{Visualizzatore DICOM 3D}} {2022}
         \resumeItemListStart
           \resumeItem{Un moderno visualizzatore di immagini mediche basato su browser per la visualizzazione di file DICOM in 2D e 3D. Gratuito e open-source.}
           \resumeItem{Sviluppato con la privacy come priorit\`a; tutta l'elaborazione avviene al 100\% lato client, senza upload su server.}
           \resumeItem{Sviluppato come strumento di apprendimento per esplorare l'analisi DICOM, il rendering volumetrico e le tecniche di visualizzazione accelerate da GPU nel browser.}
         \resumeItemListEnd
          
        \resumeProjectHeading
          {\textbf{Segmentazione Spettrale 3D}}{2021}
          \resumeItemListStart
            \resumeItem{Progetto e implemento un add-on per \href{https://www.blender.org}{Blender} per rilevare automaticamente le caratteristiche degli oggetti 3D utilizzando un algoritmo di segmentazione spettrale personalizzato.}
            \resumeItem{L'add-on ha superato le 200 scariche con numerosi feedback positivi.}
			\resumeItem{Il codice sorgente completo \`e disponibile sul mio \href{https://github.com/Rage997/Spectral-Clustering-mesh-segmentation}{profilo GitHub}.}

          \resumeItemListEnd
    \resumeSubHeadingListEnd


%
%-----------PROGRAMMING SKILLS-----------
\section{COMPETENZE}
\newcolumntype{L}{>{\raggedleft}p{0.14\textwidth}}
\newcolumntype{R}{p{0.8\textwidth}}

\begin{minipage}[ht]{0.48\textwidth}
\textbf{Linguaggi di programmazione:} \\
Python C++, C\#, OO \\ Javascript, Rust
\vspace{1em}

\textbf{Framework:} \\
Numpy, Pandas, Scikit-learn, PyTorch, Cuda, OpenMPI, \\ 
NodeJS, mongoDB, WebGL, ThreeJS
\end{minipage}
\begin{minipage}[ht]{0.48\textwidth}
\textbf{Deployment:} \\
AWS, Kubernetes, Rancher, Docker, Linux

\vspace{1em}
\textbf{Lingue:} \\
Italiano: Madrelingua \\
Inglese: Competenza Professionale Piena\\
Spagnolo: Competenza Conversazionale \\
Tedesco: Competenza Base
%Francese: Competenza professionale

\end{minipage}


%-------------------------------------------
\end{document}
